\begin{abstract}
Large language models (LLMs) can inspect source code for malicious behavior,
but they miss subtle attacks such as obfuscated droppers and typosquats whose
installation scripts look innocuous. Retrieval-augmented generation (RAG)
has been proposed as the remedy, yet prior work couples retrieval to
particular knowledge sources and mechanisms, leaving it unclear how much of
the benefit comes from retrieval itself and how much from how retrieval is
done. We isolate retrieval strategy through a controlled comparison of seven
pipelines---a zero-shot baseline and six RAG variants (dense, sparse BM25,
hybrid fusion, behavior-summary, contrastive dual-pool, and their
combination)---using the same LLM backbone, prompting, and code-only
knowledge base throughout, and evaluate on 1{,}020 test packages. We find
that code-level retrieval substantially improves detection, driven by a
large recall gain; simple sparse retrieval performs on par with heavy dense retrieval;
contrastive retrieval sharply reduces false alarms; and behavior-summary
queries are actively harmful. We also analyze the samples that no method
detects (metadata-cloned typosquats) and use a static-grounding proxy to
measure explanation faithfulness, showing that retrieval reduces unsupported
behavior claims.
\end{abstract}
\endinput
